\documentclass[11pt]{article}

\usepackage[utf8]{inputenc}
\usepackage[ngerman]{babel}
\usepackage{csquotes}

\usepackage{fullpage}
\usepackage{setspace}
\usepackage{parskip}
\usepackage{titlesec}
\usepackage{mathptmx}

\usepackage{graphicx}
\usepackage[section]{placeins}

\PassOptionsToPackage{
style=apa,
doi=false,
isbn=false,
eprint=false
}{biblatex}
\usepackage[backend=biber]{biblatex}
\addbibresource{library.bib}

\PassOptionsToPackage{hyphens}{url}

\makeatletter

\renewenvironment{abstract}
{{\bfseries\noindent{\abstractname}\par\nobreak}\footnotesize}
{\bigskip}

\renewenvironment{quote}
  {\begin{tabular}{|p{13cm}}}
  {\end{tabular}}

\titlespacing{\section}{0pt}{*3}{*1}
\titlespacing{\subsection}{0pt}{*2}{*0.5}
\titlespacing{\subsubsection}{0pt}{*1.5}{0pt}

\usepackage{authblk}

\usepackage[colorlinks = false]{hyperref}

\begin{document}

\begin{titlepage}
   \begin{center}
       \vspace*{1cm}

       \Huge
       Skribbl.AI Projektdokumentation (IQ vs. KI)
       \vspace{2.0cm}

       \includegraphics[width=0.4\textwidth]{logo.jpg}

       \vspace{1.5cm}
       \LARGE

       Josefine Sophie Busch, Malin Dulkies, Rachel Escueta, Tim-Niklas Heise, Ninoslav Kjireski, Anastasia Litvina, Anh Quang Vu-Tuyen

       \vfill

       Projektarbeit \\
       Prof. Dr. Klaus Jung\\
       Wintersemester 2018/2019\\

       \vspace{0.8cm}

       Hochschule für Technik und Wirtschaft\\
       Berlin\\

   \end{center}
\end{titlepage}

\pagebreak
\tableofcontents
\pagebreak
\listoftables
\listoffigures
\pagebreak

\section{Einleitung}
\subsection{Leitfaden}

 (Was ist das Projekt, wie ist es zustande gekommen, was ist das Ziel)
 
Das Semesterprojekt ist ein Pflichtmodul des Internationale Medieninformatik Studienganges der HTW Berlin, welches im 5. Semester angesetzt ist.

In einer Gruppe von 5-8 zufällig gewählten Studierenden werden Konzepte entwickelt und dann in Anwendungen umgesetzt, dabei können die Themen, an denen die Studierenden arbeiten, aus Bereichen der Wirtschaft oder aus einem der verschiedenen Forschungsbereiche des Studienganges kommen. 
Die Studierenden werden dabei von einer Lehrkraft bzw. einer verantwortlichen Person aus dem wirtschaftlichen Unternehmen betreut.

Das Projekt ist in drei Bereiche unterteilt: 1. der Themenvergabe/Konzeptausarbeitung, 2. dem Projektmanagement und 3. der Durchführung, dazu jedoch unter 1.2 Aufbau des Projektes mehr.

Die IMI Showtime Messe, auf welcher die Gruppen ihre Ergebnisse präsentieren, bildet den Abschluss des Semesterprojektes, dabei wird 70 Prozent (?) der Note Anhand von der Präsentation, dem Auftreten als Gruppe und dem Stand des Fortschrittes festgesetzt.
\subsection{Aufbau des Projektes}
 (Wir haben uns für agile, Scrum usw. entschieden)    1.2.1. Trello, Slack
    
Das Projekt ist, wie unter 1.1 Leitfaden schon erwähnt, in drei Bereiche unterteilt, welche im folgenden genauer und im Kontext unseres Projektes erläutert werden.

1. Themenvergabe/Konzeptausarbeitung
	Wir haben das Projekt KI vs IQ erhalten, welches von Prof. Jung 		betreut wurde.
	Die Vorgabe war, eine spielerische Anwendung zu entwickeln, 			in der Mensch gegen neuronales Netz antritt, welches dazu 				trainiert wurde, Strichzeichnungen zu erkennen.
	Der Mensch sollte etwas zeichnen und die AI sollte dieses 				erkennen und das ganze als Spiel.
	Wir entschieden uns dafür eine Website zu bauen, auf welcher 			unser Spiel dann läuft.
	Refezent auf Kapitel 2?
	
2. Projektmanagement
	Die Belegung des Semesterprojektes erfolgt mit einer 						automatischen Belegung des Projektmanagement Moduls, 				welches sich auf das agile Arbeiten mit Scrum in Software 				Projekten beschäftigt. Dem entsprechen und der Tatsache 				zugrundeliegend, dass einige Gruppenmitglieder bereits 					Erfahrung mit Scrum hatten, haben wir uns dafür entschieden, im 	Projekt unter(?) Scrum zu arbeiten bzw. einer unserer 						Bedürfnisse und Gegebenheiten angepassten Version von 				Scrum.
	Als unterstützenden Tools zum managen des Projektes haben wir 	Trello genutzt und zum managen der Gruppe bzw. als 						Kommunikations-Tool Slack.

3. Durchführung
	Die Durchführungsphase setzt sich dann konkret mit der 					Entwicklung der Anwendung zusammen.
	Wir haben HTML, CSS und Javascript für das entwickeln der 			Webumgebung und des Spieles genutzt und Python für das 				entwickeln des neuronalen Netzwerkes.
	Mehr dazu im Kapitel 3 Anforderungen.

\section{Grundlagen}
\subsection{Design}
\subsubsection{Findung Und SketchR}

Wissend, dass wir ein Spiel kreieren wollten, welches sowohl Bilderkennung durch ein neuronales Netzwerk garantiert als auch den Spielspaß für den Spieler, blieben wir bevor weitere Überlegungen angefangen werden konnten, direkt an Namen sowie Konzept des Projekts hängen. 
Abgeleitet durch das englische Wort 'to scribble' mit ein paar wenigen Änderungen um den Namen für das deutsche Publikum besser aussprechbar zu machen, wurde 'Skribbl' geboren. Das Kürzel '.AI' hingegen war für uns als Team sehr offensichtlich. Nicht nur wäre es eine wundervolle URL, wenn jemand das Projekt veröffentlichen wollen würde, sondern sie zeigt ebenfalls die nahe Verbindung zum neuronalen Netzwerk.

Folgend kam die Entscheidung dahingehend, was für eine Art von Spiel wir erschaffen wollten. Eine Zeichnung sollte erkannt werden, aber die Möglichkeiten dies umzusetzen waren zu umfangreich und so mussten wir uns auf ein Konzept einigen. Eines, welches nicht zu \textit{viel} beinhaltete und den Rahmen nicht sprengen würde, allerdings auch nicht zu einfach ausfallen würde. 

Letztendlich entschieden wir uns dafür erst einmal zu recherchieren, ob es nicht vielleicht bereits jemanden gab der ein ähnliches Konzept umgewandelt hat. Spiel oder nicht war dabei egal. 
Dabei stießen wir auf zwei Seiten verschiedene Programme (SketchR und Quick, Draw! [Google]). Geholfen haben diese bei unserer Entscheidung in welche Richtung das Spiel sich entwickeln sollte definitiv. Nach einigen Diskussionen entschieden wir uns schließlich für ein Zeichen-Spiel mit integriertem Zeitfaktor, bei dem die Künstliche Intelligenz das Bild mit Hilfe eines Treppchen Systems (ist das gesuchte Wort wahrscheinlicher als alle anderen, somit ist es auf Platz 1), erkennen sollte. 

\subsubsection{Mobile First}

Ein solches Projekt, mit Spiel-Charakter, wie man sie aus jedem App-Store kennt; kurze Level für ein Kurzweiliges Spielerlebnis machten eines schnell klar: Unsere Priorität lag auf Mobile-First. 
Dabei musste der Touchscreen, die unterschiedlichen Bildschirmgrößen sowie Funktionen und Hürden der verschiedenen Browser und Betriebssysteme mit eingeplant werden. Dennoch war es für uns keine schwere Entscheidung den Desktop vorerst außen vor zu lassen und sicher zu stellen, dass unser Spiel wenigstens auf einem Smartphone spielbar ist. 

\subsubsection{Accessibility}

Accessibility war ein Punkt, der zwar erst ein wenig später in unserem Projekt auftauchte, aber dennoch nicht von kleiner Bedeutung war. Mit dem sehr visuellen Thema des Zeichnens und der heutigen Einstellung zu Accessibility Themen, war uns schnell klar, dass wir unser Spiel so zugänglich wie möglich für alle Nutzer machen mussten. 

Ein Styleguide (Styleguide anhängen?) wurde erstellt, in welchem jede Farbe und Schriftgröße so getestet wurde, dass sie den Werten der W3C Verordnung entspricht.
Gleichzeitig wurden noch andere Details beachtet, wie zum Beispiel welche Coding-Regeln einzuhalten waren.


\subsection{Frontend}
\subsection{Neuronales Netzwerk}
\subsubsection{Dataset}
\section{Anforderungen}
Genauere Productvision, Detailierter
\section{Umsetzung, Implementierung, Ausarbeitung}
\section{Spielbeschreibung, Ziel}
\section{Bewertung}
vgl. Anforderungen mit Ergebnis

\printbibliography
\end{document}
